%% kapa: exam=1 round=36 year=2025 period=2 form=A subject=accounting subject_ko=회계학 seq=2 number=42
%% source: Q-Net 감정평가사 기출문제 (cst00309) artlSeq=5231525 file=raw/1008/5231525/2025년 제36회 감정평가사 1차시험 2교시 문제지 원본.pdf
%% conversion: kordoc 4.24.51 -> tools/gen1.py (auto); review=proofread-sonnet; flags=none
\begin{question}{42}
재무제표 표시에 관한 설명으로 옳지 않은 것은?
\begin{choices}
  \item 기업이 재무상태표에 유동자산과 비유동자산, 그리고 유동부채와 비유동부채로 구분하여 표시하는 경우, 이연법인세부채는 유동부채로 분류한다.
  \item 기업이 명확히 식별 가능한 영업주기 내에서 재화나 용역을 제공하는 경우, 재무상태표에 유동자산과 비유동자산 및 유동부채와 비유동부채를 구분하여 표시한다.
  \item 기업이 보고기간말 현재 기존의 대출계약조건에 따라 보고기간 후 적어도 12개월 이상 부채를 연장할 권리가 있다면, 보고기간 후 12개월 이내에 만기가 도래한다 하더라도 비유동부채로 분류한다.
  \item 영업주기는 영업활동을 위한 자산의 취득시점부터 그 자산이 현금이나 현금성자산으로 실현되는 시점까지 소요되는 기간을 의미하며, 정상영업주기를 명확히 식별할 수 없는 경우에는 그 기간이 12개월인 것으로 가정한다.
  \item 매입채무 그리고 종업원 및 그 밖의 영업원가에 대한 미지급비용과 같은 유동부채는 기업의 정상영업주기 내에 사용되는 운전자본의 일부이므로, 이러한 항목은 보고기간 후 12개월 후에 결제일이 도래한다 하더라도 유동부채로 분류한다.
\end{choices}
\end{question}
\answer{1}
